\documentclass[10pt,a4paper]{amsart}
\usepackage[margin=19mm]{geometry}
\usepackage{amsmath,amssymb,mathtools}
\usepackage[scr=rsfs,cal=euler]{mathalfa}
\usepackage{xfrac}
\usepackage[unicode,hidelinks,bookmarks=false]{hyperref}
\newcommand{\ie}{i.e.\ }
\newcommand{\cf}{cf.\ }
\newcommand{\resp}{respectively\ }
\newcommand{\Subsection}{\subsection}
\providecommand{\nobreakdash}{\nobreak\hbox{-}}
\setlength{\emergencystretch}{2em}
\multlinegap0pt
\usepackage{mathrsfs,color,times,textcomp,verbatim}
\hypersetup{urlcolor=blue, citecolor=red, linkcolor=blue}
\usepackage{pgfplots}
\usepackage{tikz}
\numberwithin{equation}{section}
\newtheorem{theorem}{Theorem}[section]
\newtheorem{proposition}[theorem]{Proposition}
\newtheorem{lemma}[theorem]{Lemma}
\newtheorem{conjecture}{Conjecture}
\newtheorem*{question}{Question}
\newtheorem{corollary}[theorem]{Corollary}
\newtheorem{claim}{Claim}
\allowdisplaybreaks[4]
\newtheorem{definition}[theorem]{Definition}
\newtheorem*{notation}{\textbf{Notation}}
\newtheorem*{convention}{\textbf{Convention}}
\newtheorem{remark}[theorem]{Remark}
\begin{document}
\title[Classification of finite energy positive solutions to Yamabe]{Classification of finite energy positive solutions to Yamabe-type equation on the fifteen dimensional octonionic Heisenberg group}
\author{Daowen Lin}
\date{}
\maketitle
\noindent\emph{Source and licence.} Classification of finite energy positive solutions to Yamabe-type equation on the fifteen dimensional octonionic Heisenberg group. Version arXiv:2609.17072v1 (CC BY 4.0). This material is distributed under the Creative Commons Attribution 4.0 International licence, http://creativecommons.org/licenses/by/4.0/, as recorded in the arXiv OAI licence metadata for this exact version and shown on the abstract page https://arxiv.org/abs/2609.17072v1. Full rights notice: licenses/arxiv-2609.17072-CC-BY-4.0.txt.\par\smallskip
\noindent\textit{Editorial scope.} Counted unit: the original \texttt{theorem} environment \texttt{jl} at source lines 575--584 together with its complete proof at lines 585--636; the delivered document keeps source lines 65--656 so that every referenced label and every citation resolves inside it. Native editable source is the paper's own concatenated TeX; the AMS wrapper is editorial formatting only. No publisher class, upstream script or unrelated artwork is redistributed.
\begin{equation}\label{1.1}
    \begin{aligned}
        &\Delta_G u+160u^\frac{Q+2}{Q-2}=0\\& u\in \mathscr{D}^{1,2}(G),\,\,u> 0   
    \end{aligned}
\end{equation}
where $Q=22$ is the homogeneous dimension of $G$. 

Equation (\ref{1.1}) has been studied due to its  connection with  Sobolev type inequality (\ref{1.3}) and   Yamabe type problem.

For the Heisenberg group case, which is related to the CR Yamabe problem: \textit{Given a compact, strict pseudo-convex CR manifold, find a choice of contact form for which the Webster-Tanaka pseudo-hermitian scalar curvature is constant}, Jerison and Lee \cite{JL1988} classified all finite energy positive solutions to  Yamabe type equation by introducing three family of differential identities. For further discussions,  see \cite{G2001a} \cite{G2001b} \cite{JL1987} \cite{JL1989} \cite{W2015}. Without finite energy condition, $n=1$ was addressed in \cite{CLMR2023}, with partial results for $n\geq 2$ in \cite{CLMR2023} and \cite{FV2023}, as well as in \cite{CMRW2024}. For subcritical case, see \cite{MO2023} \cite{MOW2023}. 

Ivanov, Minchev and Vassilev [\cite{IMV2010} $(n=1)$, \cite{IMV2023} $(n\geq 2)$] classified finite energy positive solutions to the Yamabe type equation on the quaternionic Heisenberg group  using differential identities. They reduced the problem to finding a positive-definite  matrix of size $7\times 7$. For subcritical case, see \cite{LZ}.

For general Heisenberg type group, Garofalo and Vassilev \cite{GV2001} made a conjecture regarding the classification of finite energy positive solutions to the  Yamabe type equation. In the same paper, the conjecture  was verified for all groups of Iwasawa type under the assumption of partial symmetry of the solution. In this paper, we confirm the conjecture of Garofalo and Vassilev on the $15$ dimensional octonionic Heisenberg group.



\begin{theorem}If $u(x,t)$ is a solution of (\ref{1.1}), then up to dilation and translation,     \begin{align}\label{1.2}
     u(x,t)=[(1+|x|^2)^2+|t|^2]^{-5}.
     \end{align}\end{theorem}







 \begin{corollary}
    Every  extremal of the following Sobolev inequality on $15$ dimensional octonionic Heisenberg group must be  (\ref{1.2}), up to dilation and translation. 
 \end{corollary}
 
[Folland and Stein \cite{FS1974}] For $v\in \mathscr{D}^{1,2}(G)$, there is a universal constant $S$ such that


    \begin{align}\label{1.3}
         \bigg(\int_{G}v^\frac{2Q}{Q-2}dV\bigg)^\frac{Q-2}{Q}\leq S\int_{G}|\nabla_G v|^2dV
     \end{align}


 
For sharp inequalities on Heisenberg type group, one can get the sharp Folland-Stein-Sobolev inequality from Jerison and Lee's result \cite{JL1988} on Heisenberg group and from Ivanov, Minchev and Vassilev's results [\cite{IMV2010}, \cite{IMV2023}] on quaternionic Heisenberg group. Frank and Lieb \cite{FL2012} established sharp Hardy-Littlewood-Sobolev inequalities on the Heisenberg group; see also Hang and Wang \cite{HW2022}. For Hardy-Littlewood-Sobolev inequality on the other Heisenberg type groups, see [\cite{CLZ},\cite{CLZ2}]. Later, Yang \cite{Y2024} gave the optimal constant in Folland-Stein-Sobolev inequality on the Heisenberg type group inspired by the work of Frank and Lieb \cite{FL2012} and Hang and Wang \cite{HW2022}, but he did not classify the minimizers.

 The proof of our result relies on a Jerison-Lee type identity along with a  prior estimates for the solution and its derivatives. The key insight of our approach is the definition (\ref{3.14}) of $B_\nu$, which allows us to analyze Heisenberg type group with several dimensions in the second layer. The main challenge arises from the complexity of the octonion, which is not associative. 
 
\textbf{Convention}: \textit{i,j,.. stand for 0,1,..,7; $\nu,\mu,...$ stand for 1,2,...,7.}


 
 The paper is organized as follows: In Section \ref{s2}, we collect some notations on the octonionic Heisenberg group. In Section \ref{s3}, we define $\mathcal{D}$ and $B_\nu$, and then we compute a Jerison-Lee type identity. In Section \ref{s4}, we use the Jerison-Lee type identity and a prior estimates to show $\mathcal{D}=B_\nu= 0$ and then we follow the method in \cite{IMV2014} to complete our proof.

 
\section{Notations}\label{s2}
In this section, we give an introduction to the octonionic Heisenberg group. We consider
\begin{align*}
    G=\mathbb{R}^{8}\times \mathbb{R}^7
\end{align*}
with $\xi=(x,t)=(x^i, t^\nu)$ where $ i=0,1,..,7; \nu=1,2,...,7$ and with the group law: given $\xi=(x,t),\zeta=(y,s)$,
\begin{align*}
    (x,t)\circ(y,s)=(x+y,T),
\end{align*}
where 
\small{
\begin{align*}
    T^1&=t^1+s^1+2(-x^0 y^1+x^1 y^0-x^2 y^3+x^3 y^2-x^4y^5+x^5 y^4+x^6 y^7-x^7 y^6),
\end{align*}
\begin{align*}
    T^2&=t^2+s^2+2(-x^0 y^2+x^1 y^3+x^2 y^0-x^3y^1-x^4y^6-x^5 y^7+x^6 y^4+x^7 y^5),
\end{align*}
\begin{align*}
    T^3&=t^3+s^3+2(-x^0 y^3-x^1 y^2+x^2 y^1+x^3 y^0-x^4y^7+x^5 y^6-x^6 y^5+x^7 y^4),
\end{align*}
\begin{align*}
    T^4&=t^4+s^4+2(-x^0 y^4+x^1 y^5+x^2 y^6+x^3 y^7+x^4 y^0-x^5 y^1-x^6 y^2-x^7 y^3),
\end{align*}
\begin{align*}
    T^5&=t^5+s^5+2(-x^0 y^5-x^1 y^4+x^2 y^7-x^3y^6+x^4y^1+x^5 y^0+x^6 y^3-x^7 y^2),
\end{align*}
\begin{align*}
    T^6&=t^6+s^6+2(-x^0 y^6-x^1 y^7-x^2 y^4+x^3 y^5+x^4y^2-x^5 y^3+x^6 y^0+x^7 y^1),
\end{align*}
\begin{align*}
    T^7&=t^7+s^7+2(-x^0 y^7+x^1 y^6-x^2 y^5-x^3y^4+x^4y^3+x^5 y^2-x^6 y^1+x^7 y^0).
\end{align*}}

We define the norm for $\xi=(x,t)$ as follow:
\begin{align*}
    |\xi|=(|x|^4+|t|^2)^\frac{1}{4},
\end{align*}
with the associated distance function 
\begin{align*}
    d(\xi,\zeta)=|\zeta^{-1}\circ\xi|,\,\,\textrm{for}\,\,\xi,\zeta \in G,
\end{align*}
where $\zeta^{-1}$ is the inverse of $\zeta$ with respect to $\circ$, i.e. $\zeta^{-1}=-\zeta$.
We define
\begin{align*}
    B_R(\xi)=\{\zeta\in G:d(\xi, \zeta)<R\},
\end{align*}
and we denote $B_R=B_R(0)$. We have
\begin{align*}
    |B_R(\xi)|=CR^Q,
\end{align*}
where $C>0$ is a positive constant, $Q=22$ is the homogeneous dimension of $G$ and $|\cdot|$ is the Lebesgue measure.



We define the following $8$ left-invariant vector fields in $G$:
\small{\begin{align*}
&\mathscr{X}^0=\partial_{x^0}+2x^1 \partial_{t^1}+2x^2 \partial_{t^2}+2x^3 \partial_{t^3}+2x^4 \partial_{t^4}+2x^5 \partial_{t^5}+2x^6 \partial_{t^6}+2x^7 \partial_{t^7},\\
&\mathscr{X}^1 =\partial_{x^1}-2x^0 \partial_{t^1}-2x^3 \partial_{t^2}+2x^2 \partial_{t^3}-2x^5 \partial_{t^4}+2x^4 \partial_{t^5}+2x^7 \partial_{t^6}-2x^6  \partial_{t^7},\\
&\mathscr{X}^2=\partial_{x^2}+2x^3 \partial_{t^1}-2x^0 \partial_{t^2}-2x^1 \partial_{t^3}-2x^6\partial_{t^4}-2x^7\partial_{t^5}+2x^4\partial_{t^6}+2x^5 \partial_{t^7},\\
&\mathscr{X}^3=\partial_{x^3}-2x^2\partial_{t^1}+2x^1 \partial_{t^2}-2x^0 \partial_{t^3}-2x^7 \partial_{t^4}+2x^6 \partial_{t^5}-2x^5 \partial_{t^6}+2x^4 \partial_{t^7},\\
&\mathscr{X}^4=\partial_{x^4}+2x^5 \partial_{t^1}+2x^6 \partial_{t^2}+2x^7 \partial_{t^3}-2x^0\partial_{t^4}-2x^1 \partial_{t^5}-2x^2 \partial_{t^6}-2x^3\partial_{t^7},\\
&\mathscr{X}^5=\partial_{x^5}-2x^4 \partial_{t^1}+2x^7 \partial_{t^2}-2x^6 \partial_{t^3}+2x^1 \partial_{t^4}-2x^0\partial_{t^5}+2x^3\partial_{t^6}-2x^2 \partial_{t^7},\\
&\mathscr{X}^6=\partial_{x^6}-2x^7 \partial_{t^1}-2x^4\partial_{t^2}+2x^5\partial_{t^3}+2x^2 \partial_{t^4}-2x^3 \partial_{t^5}-2x^0\partial_{t^6}+2x^1 \partial_{t^7},\\
&\mathscr{X}^7=\partial_{x^7}+2x^6 \partial_{t^1}-2x^5 \partial_{t^2}-2x^4 \partial_{t^3}+2x^3 \partial_{t^4}+2x^2 \partial_{t^5}-2x^1 \partial_{t^6}-2x^0 \partial_{t^7}.   
 \end{align*}}
In fact, the above vector fields are the kernel of the following $7$ contact forms,
\small{\begin{align*}
    \eta_1=&dt^1+2(-x^1 dx^0+x^0 dx^1-x^3 dx^2+x^2 dx^3-x^5dx^4+x^4 dx^5+x^7 dx^6-x^6 dx^7),\\
    \eta_2=&dt^2+2(-x^2 dx^0+x^3 dx^1+x^0 dx^2-x^1dx^3-x^6 dx^4-x^7 dx^5+x^4 dx^6+x^5 dx^7),\\
    \eta_3=&dt^3+2(-x^3dx^0-x^2 dx^1+x^1dx^2+x^0 dx^3-x^7 dx^4+x^6dx^5-x^5 dx^6+x^4 dx^7),\\
    \eta_4=&dt^4+2(-x^4 dx^0+x^5dx^1+x^6 dx^2+x^7 dx^3+x^0 dx^4-x^1dx^5-x^2dx^6-x^3 dx^7),\\
    \eta_5=&dt^5+2(-x^5 dx^0-x^4dx^1+x^7 dx^2-x^6 dx^3+x^1 dx^4+x^0 dx^5+x^3 dx^6-x^2 dx^7),\\
    \eta_6=&dt^6+2(-x^6 dx^0-x^7 dx^1-x^4dx^2+x^5 dx^3+x^2dx^4-x^3 dx^5+x^0 dx^6+x^1 dx^7),\\
    \eta_7=&dt^7+2(-x^7 dx^0+x^6 dx^1-x^5 dx^2-x^4 dx^3+x^3 dx^4+x^2 dx^5-x^1 dx^6+x^0 dx^7).
    \end{align*}}
The following forms $\omega_\nu$ are associated to $d\eta_\nu$,
\small{\begin{align*}
    \omega_1=& dx^0\wedge dx^1+dx^2\wedge dx^3+dx^4\wedge dx^5+dx^7\wedge dx^6,\\
    \omega_2=&dx^0\wedge dx^2+dx^3\wedge dx^1+dx^4\wedge dx^6+dx^5\wedge dx^7,\\
    \omega_3=&dx^0\wedge dx^3+dx^1\wedge dx^2+dx^4\wedge dx^7+dx^6\wedge dx^5,\\
    \omega_4=&dx^0\wedge dx^4+dx^5\wedge dx^1+dx^6\wedge dx^2+dx^7\wedge dx^3,\\
    \omega_5=&dx^0\wedge dx^5+dx^1\wedge dx^4+dx^7\wedge dx^2+dx^3\wedge dx^6,
    \\\omega_6=&dx^0\wedge dx^6+dx^1\wedge dx^7+dx^2\wedge dx^4+dx^5\wedge dx^3,
    \\  
    \omega_7=&dx^0\wedge dx^7+dx^6\wedge dx^1+dx^2\wedge dx^5+dx^3\wedge dx^4.   \end{align*}}

We define some operators related to the octonionic multiplication for convenience:

   
\begin{equation*}
   \left(
    \begin{array}{ccccccccc}
&\mathscr{X}^0&\mathscr{X}^1&\mathscr{X}^2& \mathscr{X}^3& \mathscr{X}^4&\mathscr{X}^5& \mathscr{X}^6& \mathscr{X}^7\\ 
I_1&\mathscr{X}^1&-\mathscr{X}^0&\mathscr{X}^3& -\mathscr{X}^2& \mathscr{X}^5&-\mathscr{X}^4& -\mathscr{X}^7& \mathscr{X}^6\\ 
I_2&\mathscr{X}^2&-\mathscr{X}^3&-\mathscr{X}^0& \mathscr{X}^1& \mathscr{X}^6&\mathscr{X}^7& -\mathscr{X}^4& -\mathscr{X}^5\\ 
I_3&\mathscr{X}^3&\mathscr{X}^2&-\mathscr{X}^1& -\mathscr{X}^0& \mathscr{X}^7&-\mathscr{X}^6& \mathscr{X}^5& -\mathscr{X}^4\\ 
I_4&\mathscr{X}^4&-\mathscr{X}^5&-\mathscr{X}^6& -\mathscr{X}^7& -\mathscr{X}^0&\mathscr{X}^1& \mathscr{X}^2& \mathscr{X}^3\\
I_5&\mathscr{X}^5&\mathscr{X}^4&-\mathscr{X}^7& \mathscr{X}^6& -\mathscr{X}^1&-\mathscr{X}^0& -\mathscr{X}^3& \mathscr{X}^2\\ 
I_6&\mathscr{X}^6&\mathscr{X}^7&\mathscr{X}^4& -\mathscr{X}^5& -\mathscr{X}^2&\mathscr{X}^3& -\mathscr{X}^0& -\mathscr{X}^1\\ 
I_7&\mathscr{X}^7&-\mathscr{X}^6&\mathscr{X}^5& \mathscr{X}^4& -\mathscr{X}^3&-\mathscr{X}^2& \mathscr{X}^1& -\mathscr{X}^0\\ 
\end{array}
    \right)
    \end{equation*}
For example, $I_5\mathscr{X}^2=I_5I_2\mathscr{X}^0=-I_7\mathscr{X}^0=-\mathscr{X}^7$. 

We list the Lie brackets of the horizontal vector fields.
\begin{lemma}\label{lemma2.1}
We have
  \begin{equation*}
      \begin{aligned}
\big[\mathscr{X}^0,\,\mathscr{X}^1\big]=\big[\mathscr{X}^2,\,\mathscr{X}^3\big]=\big[\mathscr{X}^4,\,\mathscr{X}^5\big]=\big[\mathscr{X}^7,\,\mathscr{X}^6\big]=-4\partial_{t^1},
      \end{aligned}
  \end{equation*}
  \begin{equation*}
      \begin{aligned}
\big[\mathscr{X}^0,\,\mathscr{X}^2\big]=\big[\mathscr{X}^3,\,\mathscr{X}^1\big]=\big[\mathscr{X}^4,\,\mathscr{X}^6\big]=\big[\mathscr{X}^5,\,\mathscr{X}^7\big]=-4\partial_{t^2},
      \end{aligned}
  \end{equation*}
  \begin{equation*}
      \begin{aligned}
\big[\mathscr{X}^0,\,\mathscr{X}^3\big]=\big[\mathscr{X}^1,\,\mathscr{X}^2\big]=\big[\mathscr{X}^4,\,\mathscr{X}^7\big]=\big[\mathscr{X}^6,\,\mathscr{X}^5\big]=-4\partial_{t^3},
      \end{aligned}
  \end{equation*}
  \begin{equation*}
      \begin{aligned}
\big[\mathscr{X}^0,\,\mathscr{X}^4\big]=\big[\mathscr{X}^5,\,\mathscr{X}^1\big]=\big[\mathscr{X}^6,\,\mathscr{X}^2\big]=\big[\mathscr{X}^7,\,\mathscr{X}^3\big]=-4\partial_{t^4},
      \end{aligned}
  \end{equation*}
  \begin{equation*}
      \begin{aligned}
\big[\mathscr{X}^0,\,\mathscr{X}^5\big]=\big[\mathscr{X}^1,\,\mathscr{X}^4\big]=\big[\mathscr{X}^7,\,\mathscr{X}^2\big]=\big[\mathscr{X}^3,\,\mathscr{X}^6\big]=-4\partial_{t^5},
      \end{aligned}
  \end{equation*}
  \begin{equation*}
      \begin{aligned}
\big[\mathscr{X}^0,\,\mathscr{X}^6\big]=\big[\mathscr{X}^1,\,\mathscr{X}^7\big]=\big[\mathscr{X}^2,\,\mathscr{X}^4\big]=\big[\mathscr{X}^5,\,\mathscr{X}^3\big]=-4\partial_{t^6},
      \end{aligned}
  \end{equation*}
  \begin{equation*}
\begin{aligned}
\big[\mathscr{X}^0,\,\mathscr{X}^7\big]=\big[\mathscr{X}^6,\,\mathscr{X}^1\big]=\big[\mathscr{X}^2,\,\mathscr{X}^5\big]=\big[\mathscr{X}^3,\,\mathscr{X}^4\big]=-4\partial_{t^7}.
      \end{aligned}
  \end{equation*}
\end{lemma}

We should be careful since octonionic multiplication is not associative. In general, we don't have $I_\nu I_\mu\mathscr{X}^i=(I_\nu I_\mu)\mathscr{X}^i$. However, we still have the following relations.
\begin{lemma}\label{lemma2.2}
For $i=0,1,...,7$ and $\nu,\mu=1,2,..,7$, we have
    \begin{equation*}
        \begin{aligned}
            I_\nu I_\mu I_\nu \mathscr{X}^i=I_\mu \mathscr{X}^i,\,\,\textrm{for}\,\nu\neq \mu;I_\nu I_\mu I_\nu \mathscr{X}^i=-I_\mu \mathscr{X}^i,\,\,\textrm{for}\,\nu= \mu.        \end{aligned}
    \end{equation*}
\end{lemma}
\begin{lemma}\label{lemma2.3}
    For $\nu,\mu=1,2,...,7$, $\nu\neq \mu$ we have
    \begin{align*}
        \Sigma_{i=0}^7\mathscr{X}^i h I_\nu\mathscr{X}^i h=0,\Sigma_{i=0}^7 I_\mu \mathscr{X}^ih I_\nu\mathscr{X}^i h=0,\Sigma_{i=0}^7 (I_\mu \mathscr{X}^i I_\nu\mathscr{X}^i h+I_\nu\mathscr{X}^iI_\mu\mathscr{X}^ih)=0.    \end{align*}
\end{lemma}


The sub-Laplacian operator on the octonionic Heisenberg group defined as
\begin{align*}
    \Delta_G h:=\Sigma_{j=0}^7(\mathscr{X}^j)^2h.
\end{align*}
And 
\begin{align*}
    |\nabla_G h|^2:=\Sigma_{i=0}^7(\mathscr{X}^j h)^2=\Sigma_{i=0}^7(I_\nu\mathscr{X}^i h)^2,
    \end{align*}
for $\nu=1,2,...,7.$



    
\section{Jerison-Lee type identity }\label{s3} 
In this section, we derive a Jerison-Lee type identity.


Suppose $u>0$ is a solution of (\ref{1.1}), we set  $u=h^{-\frac{Q-2}{4}}$, then $h$ satisfies
\begin{align}\label{3.1}
    \Delta_G h=6h^{-1}|\nabla_G h|^2+32.
\end{align}
The following function is very important in our proof,
\begin{align}\label{3.2}
    f=4+\frac{|\nabla_G h|^2}{4h}.
    \end{align}

We define $\mathcal{E}(\mathscr{X}^i,\mathscr{X}^j)=0$  and $\mathcal{D}(\mathscr{X}^i,\mathscr{X}^j)$, for $i,j =0,1,...,7.$
\begin{align*}
   \mathcal{E}(\mathscr{X}^i,\mathscr{X}^j)=&\frac{h^{-1}}{8}\bigg[\mathscr{X}^j\mathscr{X}^ih+\Sigma_{\nu=1}^7(I_\nu \mathscr{X}^j)(I_\nu \mathscr{X}^i) h-\Delta_G h \delta_{ij}+8\Sigma_{\nu=1}^7h_{t^\nu}\omega_\nu(\mathscr{X}^i,\mathscr{X}^j)\bigg]\\&-\frac{4h^{-2}}{8}\bigg[\mathscr{X}^jh\mathscr{X}^ih+\Sigma_{\nu=1}^7(I_\nu \mathscr{X}^j)h(I_\nu \mathscr{X}^i) h-|\nabla_G h|^2\delta_{ij}\bigg]\\
   =&\frac{h^{-1}}{8}\bigg[\mathscr{X}^j\mathscr{X}^ih+\Sigma_{\nu=1}^7(I_\nu \mathscr{X}^j)(I_\nu \mathscr{X}^i) h\bigg]\\&-\frac{4h^{-2}}{8}\bigg[\mathscr{X}^jh\mathscr{X}^ih+\Sigma_{\nu=1}^7(I_\nu \mathscr{X}^j)h(I_\nu \mathscr{X}^i) h\bigg]\\&
   -\frac{h^{-1}}{8}(32+2h^{-1}|\nabla_G h|^2)\delta_{ij}+h^{-1}\Sigma_{\nu=1}^7h_{t^\nu}\omega_\nu(\mathscr{X}^i,\mathscr{X}^j)\\=&0,\\
   \mathcal{D}(\mathscr{X}^i,\mathscr{X}^j)=&\frac{h^{-1}}{8}\bigg[7\mathscr{X}^j\mathscr{X}^ih-\Sigma_{\nu=1}^7(I_\nu \mathscr{X}^j)(I_\nu \mathscr{X}^i) h-24\Sigma_{\nu=1}^7h_{t^\nu}\omega_\nu(\mathscr{X}^i,\mathscr{X}^j)\bigg].
    \end{align*}
Then from Lemma \ref{lemma2.1} we have
\begin{equation}\label{3.3}
\begin{aligned}
     &\mathcal{E}(\mathscr{X}^i,\mathscr{X}^j)=\mathcal{E}(\mathscr{X}^j,\mathscr{X}^i)=0,\mathcal{D}(\mathscr{X}^i,\mathscr{X}^j)=\mathcal{D}(\mathscr{X}^j,\mathscr{X}^i),\\
     &\Sigma_{i,j=0}^7\mathcal{D}(\mathscr{X}^i,\mathscr{X}^j)\delta_{ij}=0.    \end{aligned}
\end{equation}
And for $\nu=1,2,...,7$,
\begin{align}\label{3.4}
    \mathcal{D}(\mathscr{X}^i,\mathscr{X}^j)+\Sigma_{\nu=1}^7\mathcal{D}(I_\nu \mathscr{X}^i,I_\nu \mathscr{X}^j)=0.
\end{align}
We denote 
    \begin{align}
        \mathcal{D}^2=\Sigma_{i,j=0}^7\mathcal{D}(\mathscr{X}^i,\mathscr{X}^j)^2.   \end{align}
We also have for $\nu,\mu=1,2,...,7$ from Lemma \ref{lemma2.1},
 \begin{equation}\label{3.6}
 \begin{aligned}
& \Sigma_{i,j=0}^7\mathcal{D}(\mathscr{X}^i,I_\nu \mathscr{X}^j)\omega_\mu(\mathscr{X}^i,\mathscr{X}^j)=-\Sigma_{i=0}^7\mathcal{D}(I_\mu \mathscr{X}^i, I_\nu \mathscr{X}^i)=0.\\&\Sigma_{i,j=0}^7\mathcal{D}(\mathscr{X}^i, I_\nu\mathscr{X}^j)\delta_{ij}=0.
 \end{aligned}
 \end{equation}
 We define  
\begin{align*}
   E(\mathscr{X}^j)=&h^{-1}\Sigma_{i=0}^7\mathcal{E}( \mathscr{X}^j, \mathscr{X}^i)\mathscr{X}^i h   \\=&\frac{h^{-2}}{8}\bigg[\Sigma_{i=0}^7\mathscr{X}^i h\mathscr{X}^i\mathscr{X}^jh+\Sigma_{\nu=1}^7\Sigma_{i=0}^7\mathscr{X}^ihI_\nu \mathscr{X}^i(I_\nu \mathscr{X}^j)h\\&+(-32-6h^{-1}|\nabla_G h|^2)\mathscr{X}^j h\bigg]+h^{-2}\Sigma_{\nu=1}^7h_{t^\nu}(I_\nu \mathscr{X}^j)h\\=&0,\\
    D(\mathscr{X}^j)=&h^{-1}\Sigma_{i=0}^7\mathcal{D}( \mathscr{X}^j, \mathscr{X}^i)\mathscr{X}^i h\\=&\frac{h^{-2}}{8}\bigg[7\Sigma_{i=0}^7\mathscr{X}^i h\mathscr{X}^i\mathscr{X}^j h-\Sigma_{\nu=1}^7 \mathscr{X}^ihI_\nu \mathscr{X}^i(I_\nu \mathscr{X}^j) h\bigg]-3h^{-2}\Sigma_{\nu=1}^7h_{t^\nu}(I_\nu \mathscr{X}^j)h,\\
    F_\nu(\mathscr{X}^j)=&h^{-1}\Sigma_{i=0}^7\mathcal{D}(I_\nu \mathscr{X}^j,I_\nu \mathscr{X}^i)\mathscr{X}^i h,\,\textrm{for}\,\nu=1,2,...,7.
\end{align*}   
Then from (\ref{3.4}) we have
\begin{align}\label{3.7}
    D(\mathscr{X}^j)+\Sigma_{\nu=1}^7 F_\nu (\mathscr{X}^j)=0.
\end{align}
We define

\begin{equation}\label{3.8}
\begin{aligned}
    \mathbb{D}(\mathscr{X}^i,\mathscr{X}^j,\mathscr{X}^k)=h^{-1}\bigg[&D(\mathscr{X}^i,\mathscr{X}^k)\mathscr{X}^jh+\Sigma_{\nu=1}^7D(I_\nu \mathscr{X}^i,\mathscr{X}^k)I_\nu \mathscr{X}^j h\\&+D(\mathscr{X}^j,\mathscr{X}^k)\mathscr{X}^ih+\Sigma_{\nu=1}^7D(I_\nu \mathscr{X}^j,\mathscr{X}^k)I_\nu \mathscr{X}^i h\bigg].
\end{aligned}
\end{equation}





\begin{lemma} \label{lemma 3.1}
For $\nu=1,2,..,7$, we have
\begin{align}\label{3.9}
        \Sigma_{i,j=0}^7(I_\nu \mathscr{X}^i )\mathscr{X}^j h \mathcal{D}(I_\nu \mathscr{X}^i,I_\nu \mathscr{X}^j)= \Sigma_{i,j=0}^7 \mathscr{X}^i \mathscr{X}^j h \mathcal{D}( \mathscr{X}^i,I_\nu \mathscr{X}^j)=0.
    \end{align}
    \end{lemma}
  \begin{proof}
    We  have
    \begin{equation*}
        \begin{aligned}
    &\Sigma_{i,j=0}^7\mathcal{D}(\mathscr{X}^j,\mathscr{X}^i)\mathcal{D}(\mathscr{X}^i,I_\nu \mathscr{X}^j)=\Sigma_{i,j=0}^7\mathcal{D}(\mathscr{X}^i,\mathscr{X}^j)\mathcal{D}(\mathscr{X}^i,I_\nu \mathscr{X}^j)\\=&\Sigma_{i,j=0}^7\mathcal{D}(\mathscr{X}^i,-I_\nu \mathscr{X}^j)\mathcal{D}(\mathscr{X}^i, \mathscr{X}^j)=-\Sigma_{i,j=0}^7\mathcal{D}(\mathscr{X}^i,I_\nu \mathscr{X}^j)\mathcal{D}(\mathscr{X}^j, \mathscr{X}^i)\\=&0.
    \end{aligned}    \end{equation*}
That is
\begin{equation}\label{3.10}
\begin{aligned}
      \Sigma_{i,j=0}^7\bigg[7\mathscr{X}^i\mathscr{X}^j h-\Sigma_{\mu=1}^7(I_\mu \mathscr{X}^i) I_\mu \mathscr{X}^j h-24\Sigma_{\mu=1}^7h_{t^\mu}\omega_\mu (\mathscr{X}^j,\mathscr{X}^i)\bigg]D(\mathscr{X}^i,I_\nu \mathscr{X}^j)=0. 
    \end{aligned}
\end{equation}
We also have
\begin{align}
\mathscr{X}^i\mathscr{X}^jh+\Sigma_{\mu=1}^7(I_\mu \mathscr{X}^j)(I_\mu \mathscr{X}^i)h-\Delta_G h \delta_{ij}+8\Sigma_{\mu=1}^7h_{t^\mu}\omega_\mu (\mathscr{X}^i,\mathscr{X}^j)=0.
\end{align}
Last, from (\ref{3.6}), we have 
\begin{align}
   \Sigma_{\mu=1}^7 \Sigma_{i,j=0}^7 \mathcal{D}(\mathscr{X}^i,I_\nu \mathscr{X}^j)h_{t^\mu}\omega_\mu (\mathscr{X}^i,\mathscr{X}^j)=0, \Sigma_{i,j=0}^7\mathcal{D}(\mathscr{X}^i, I_\nu\mathscr{X}^j)\delta_{ij}=0.\end{align}
Combining these, we get (\ref{3.9}).
 \end{proof}  

  We compute the derivative of $f$ defined in (\ref{3.2}).
  \begin{lemma}\label{lemma3.2} For $i=0,1,...,7,$ we have
  \begin{equation}\label{3.13}       \begin{aligned}
    2\mathscr{X}^i f&=h(D(\mathscr{X}^i)+E(\mathscr{X}^i))- 2\Sigma_{\nu=1}^7h^{-1}h_{t^\nu}(I_\nu \mathscr{X}^i)h+h^{-1}f\mathscr{X}^ih\\&=hD(\mathscr{X}^i)- 2\Sigma_{\nu=1}^7h^{-1}h_{t^\nu}(I_\nu \mathscr{X}^i)h+h^{-1}f\mathscr{X}^ih.\end{aligned}   \end{equation}
          \end{lemma}
\begin{proof}

 First, we have
    \begin{align*}
    D(\mathscr{X}^i)+E(\mathscr{X}^i)=&h^{-2}\Sigma_{j=0}^7\mathscr{X}^jh\mathscr{X}^j\mathscr{X}^ih-2\Sigma_{\nu=1}^7h^{-2}h_{t^\nu}(I_\nu \mathscr{X}^i)h\\&-\frac{3}{4}h^{-3}|\nabla_G h|^2\mathscr{X}^ih-4h^{-2}\mathscr{X}^ih.
\end{align*}
Second,
\begin{align*}
    \Sigma_{j=0}^7\mathscr{X}^j h \mathscr{X}^j\mathscr{X}^ih=\Sigma_{j=0}^7\mathscr{X}^j h \mathscr{X}^i\mathscr{X}^j h+4\Sigma_{\nu=1}^7h_{t^\nu}(I_\nu \mathscr{X}^i) h.
    \end{align*} 
    Last, we have
\begin{align*}
    2\mathscr{X}^if&=-\frac{1}{2}h^{-2}|\nabla_G h|^2\mathscr{X}^ih+h^{-1}\Sigma_{j=0}^7\mathscr{X}^j h \mathscr{X}^i\mathscr{X}^j h\\&=-2h^{-1}f\mathscr{X}^ih+8h^{-1}\mathscr{X}^ih+h^{-1}\Sigma_{j=0}^7\mathscr{X}^jh\mathscr{X}^j\mathscr{X}^ih-4\Sigma_{\nu=1}^7h^{-1}h_{t^\nu}(I_\nu \mathscr{X}^i)h.
    \end{align*}
    Combing these, we have (\ref{3.13}).
    \end{proof}
For $i=0,1,...,7; \nu=1,2,...,7$, we define $B_\nu(\mathscr{X}^i)$ as follow:
\begin{align}\label{3.14}
    B_\nu(\mathscr{X}^i)=2h^{-1}I_\nu \mathscr{X}^i h_{t^\nu}-\Sigma_{\mu=1}^7h^{-2}h_{t^\mu}(I_\mu \mathscr{X}^i)h+\frac{1}{2}h^{-2}f\mathscr{X}^ih.
\end{align}
 We also define
\begin{align}\label{3.15}
    B(\mathscr{X}^i):=\Sigma_{\nu=1}^7 B_\nu (\mathscr{X}^i)=2h^{-1}\Sigma_{\nu=1}^7I_\nu \mathscr{X}^i h_{t^\nu}-7\Sigma_{\mu=1}^7h^{-2}h_{t^\mu}(I_\mu \mathscr{X}^i)h+\frac{7}{2}h^{-2}f\mathscr{X}^ih.\end{align}

We discuss the relation between $\mathcal{D}$ and $B_\nu$.
\begin{lemma}\label{lemma3.3} For $j=0,1,...,7$, we have
\begin{equation}\label{3.16}
     \begin{aligned}
h^{11}\Sigma_{i=0}^7\mathscr{X}^i \bigg[h^{-11}\mathcal{D}(\mathscr{X}^j,\mathscr{X}^i)\bigg]=&3\bigg[\frac{7E(\mathscr{X}^j)-D(\mathscr{X}^j)}{2}+B(\mathscr{X}^j)\bigg]\\=&-\frac{3}{2}D(\mathscr{X}^j)+3B(\mathscr{X}^j).
        \end{aligned}\end{equation}
   \end{lemma}
\begin{proof}
We start with
    \begin{align*}
    &h^{11}\Sigma_{i=0}^7\mathscr{X}^i\bigg[h^{-11}\mathcal{D}(\mathscr{X}^j,\mathscr{X}^i)\bigg]\\=&\frac{h^{-1}}{8}\bigg(7\Sigma_{i=0}^7   
    (\mathscr{X}^i)^2\mathscr{X}^jh-        
    \Sigma_{i=0}^7    
    \Sigma_{\nu=1}^7\mathscr{X}^i(I_\nu\mathscr{X}^i)(I_\nu \mathscr{X}^j)h-24\Sigma_{\nu=1}^7(I_\nu \mathscr{X}^j)h_{t^\nu}\bigg)-12D(\mathscr{X}^j).
\end{align*}
From Lemma \ref{lemma2.1}, we have
\begin{align*}
    \Sigma_{i=0}^7\mathscr{X}^i (I_\nu \mathscr{X}^i)=-16\partial_{t^\nu},
\end{align*}
and
\begin{align*}
    7\Sigma_{i=0}^7 (\mathscr{X}^i)^2\mathscr{X}^jh=7\Sigma_{i=0}^7\mathscr{X}^j(\mathscr{X}^i)^2 h+56 \Sigma_{\nu=1}^7(I_\nu \mathscr{X}^j)h_{t^\nu} .
    \end{align*}
    Using (\ref{3.1}) and Lemma \ref{lemma3.2}, we have
    \begin{align*}
       & 7\mathscr{X}^j(\Delta_G h)=168\mathscr{X}^jf\\=&84h(D(\mathscr{X}^j)+E(\mathscr{X}^j))-168\Sigma_{\nu=1}^7h^{-1}h_{t^\nu}(I_\nu \mathscr{X}^j)h+84h^{-1}f\mathscr{X}^jh.
    \end{align*}
  Combining these, we have (\ref{3.16}).  
    \end{proof}









   
  

 To derive Jerison-Lee type identity, we need the following propositions.        
\begin{proposition}
    \begin{equation}
    \begin{aligned}
O_1=&h^{11}\Sigma_{i,j=0}^7\mathscr{X}^i\bigg[h^{-11}f\mathscr{X}^j h\mathcal{D}(\mathscr{X}^i,\mathscr{X}^j)\bigg]\\=& fh\mathcal{D}^2+\Sigma_{i=0}^7D(\mathscr{X}^i)\bigg[\frac{1}{2}h^2D(\mathscr{X}^i)-\Sigma_{\nu=1}^7h_{t^\nu}I_\nu \mathscr{X}^i h+\frac{1}{2}f\mathscr{X}^i h\bigg]\\&+\Sigma_{i=0}^7 f\mathscr{X}^i h\bigg[-\frac{3}{2}D(\mathscr{X}^i)+3B(\mathscr{X}^i)\bigg] 
\end{aligned}
    \end{equation}
\end{proposition}
\begin{proof}
We use Lemma \ref{lemma3.2} and Lemma \ref{lemma3.3} with
    \begin{equation*}
    \begin{aligned}
&h^{11}\Sigma_{i,j=0}^7\mathscr{X}^i\bigg[h^{-11}f\mathscr{X}^j h\mathcal{D}(\mathscr{X}^i,\mathscr{X}^j)\bigg]\\=&h^{11}\Sigma_{i,j=0}^7     f\mathscr{X}^j h\mathscr{X}^i\bigg[h^{-11}\mathcal{D}(\mathscr{X}^i,\mathscr{X}^j)\bigg]+f\Sigma_{i,j=0}^7\mathscr{X}^i\mathscr{X}^j h \mathcal{D}(\mathscr{X}^i,\mathscr{X}^j)+\Sigma_{i=0}^7h\mathscr{X}^if D(\mathscr{X}^i).
\end{aligned}
    \end{equation*}
    It suffices to prove
\begin{align*}
    &h^{-1}\Sigma_{i,j=0}^7\mathscr{X}^i\mathscr{X}^j h \mathcal{D}(\mathscr{X}^i,\mathscr{X}^j)\\=&\frac{h^{-1}}{8}\Sigma_{i,j=0}^7\bigg[7\mathscr{X}^i\mathscr{X}^j h-\Sigma_{\nu=1}^7(I_\nu \mathscr{X}^i )(I_\nu \mathscr{X}^j)h-24\Sigma_{\nu=1}^7h_{t^\nu}\omega_\nu (\mathscr{X}^j,\mathscr{X}^i) \bigg]\mathcal{D}(\mathscr{X}^i,\mathscr{X}^j)    \\=&\mathcal{D}^2.
    \end{align*}    
  It follows from
  \begin{align*}
       &\frac{1}{8}h^{-1}\Sigma_{i,j=0}^7\bigg[\mathscr{X}^i\mathscr{X}^j h+
       \Sigma_{\nu=1}^7(I_\nu \mathscr{X}^i )(I_\nu \mathscr{X}^j)h       
       \bigg]\mathcal{D}(\mathscr{X}^i,\mathscr{X}^j)\\=&\frac{1}{8}h^{-1}\Sigma_{i,j=0}^7 \mathscr{X}^i\mathscr{X}^jh\bigg[\mathcal{D}(\mathscr{X}^i,\mathscr{X}^j)+\Sigma_{\nu=1}^7\mathcal{D}(I_\nu \mathscr{X}^i,I_\nu \mathscr{X}^j)    \bigg]\\=&0       
       \end{align*}
       and
       \begin{align*}
\Sigma_{\nu=1}^7\Sigma_{i,j=0}^7\mathcal{D}(\mathscr{X}^i,\mathscr{X}^j)\omega_\nu(\mathscr{X}^j,\mathscr{X}^i)h_{t^\nu}=0.
\end{align*}

\end{proof}



\begin{proposition}
    \begin{equation}\label{3.18}
    \begin{aligned}
O_2=&h^{11}\Sigma_{\nu=1}^7\Sigma_{i,j=0}^7I_\nu \mathscr{X}^i\bigg[h^{-11}h_{t^\nu}\mathscr{X}^j h\mathcal{D}(I_\nu \mathscr{X}^i,I_\nu \mathscr{X}^j)\bigg]\\=&\Sigma_{i=0}^7\bigg[\Sigma_{\nu=1}^7\frac{1}{2}h^2F_\nu(\mathscr{X}^i)B_\nu (\mathscr{X}^i)-\frac{1}{2}\Sigma_{\mu=1}^7h_{t^\mu} D(\mathscr{X}^i) (I_\mu\mathscr{X}^i)h +\frac{1}{4}fD(\mathscr{X}^i)\mathscr{X}^i h \bigg]\\&
-\Sigma_{\nu=1}^7\Sigma_{i=0}^7h_{t^\nu}I_\nu\mathscr{X}^i h\bigg[\frac{-3D( \mathscr{X}^i)}{2}+3B(\mathscr{X}^i)\bigg]\end{aligned}
    \end{equation}
\end{proposition}
\begin{proof}
Using Lemma \ref{3.1}, Lemma \ref{3.3} and (\ref{3.7}), we have
     \begin{align*}
&h^{11}\Sigma_{\nu=1}^7\Sigma_{i,j=0}^7I_\nu \mathscr{X}^i\bigg[h^{-11}h_{t^\nu}\mathscr{X}^j h\mathcal{D}(I_\nu \mathscr{X}^i,I_\nu \mathscr{X}^j)\bigg]\\=&\Sigma_{\nu=1}^7\Sigma_{i=0}^7hF_\nu(\mathscr{X}^i) \bigg[\frac{1}{2} h B_\nu (\mathscr{X}^i)+\frac{1}{2}\Sigma_{\mu=1}^7 h^{-1}h_{t^\mu} (I_\mu\mathscr{X}^i)h -\frac{1}{4}h^{-1}f\mathscr{X}^i h \bigg]\\&+\Sigma_{\nu=1}^7\Sigma_{j=0}^7h_{t^\nu}\mathscr{X}^j h\bigg[\frac{-3D(I_\nu \mathscr{X}^j)}{2}+3B(I_\nu\mathscr{X}^j)\bigg] \\=&\Sigma_{i=0}^7\bigg[\Sigma_{\nu=1}^7\frac{1}{2}h^2F_\nu(\mathscr{X}^i)B_\nu (\mathscr{X}^i)-\frac{1}{2}\Sigma_{\mu=1}^7 h_{t^\mu}D(\mathscr{X}^i) (I_\mu\mathscr{X}^i)h +\frac{1}{4}fD(\mathscr{X}^i)\mathscr{X}^i h \bigg]\\&
-\Sigma_{\nu=1}^7\Sigma_{i=0}^7h_{t^\nu}I_\nu\mathscr{X}^i h\bigg[\frac{-3D( \mathscr{X}^i)}{2}+3B(\mathscr{X}^i)\bigg]. 
\end{align*}




\end{proof}






\begin{proposition}
\begin{equation}\label{3.19}
    \begin{aligned}
       O_3=& h^{11}\Sigma_{\nu=1}^7\Sigma_{i=0}^7I_\nu\mathscr{X}^i \bigg[h^{-10}B_\nu(\mathscr{X}^i)h_{t^\nu}\bigg]   \\=&\frac{1}{2}h^2\Sigma_{\nu=1}^7B_\nu^2 +\Sigma_{i=0}^7\bigg[-\frac{1}{4}\Sigma
_{\nu=1}^7D(\mathscr{X}^i)I_\nu\mathscr{X}^i h h_{t^\nu}+\Sigma_{\nu=1}^7B(\mathscr{X}^i)I_\nu{\mathscr{X}^i}hh_{t^\nu}\\&-\frac{1}{4}fB(\mathscr{X}^i)\mathscr{X}^i h\bigg]
       \end{aligned} 
       \end{equation}\end{proposition}
    

\begin{proof}
From the definition of (\ref{3.14}), we have
\begin{equation}
    \begin{aligned}
    &\Sigma_{\nu=1}^7\Sigma_{i=0}^7h^{11}I_\nu\mathscr{X}^i \bigg[h^{-10}B_\nu(\mathscr{X}^i)h_{t^\nu}\bigg]\\=&-12\Sigma_{\nu=1}^7\Sigma_{i=0}^7h_{t^\nu}B_\nu(\mathscr{X}^i)I_\nu\mathscr{X}^i h+\Sigma_{\nu=1}^7\Sigma_{i=0}^7 hI_\nu\mathscr{X}^i h_{t^\nu}B_\nu(\mathscr{X}^i)\\&+h^{-1}\Sigma_{\nu=1}^7h_{t^\nu}\Sigma_{i=0}^7\bigg[2I_\nu\mathscr{X}^i h I_\nu\mathscr{X}^i h_{t^\nu}+2h I_\nu\mathscr{X}^i I_\nu\mathscr{X}^i h_{t^\nu}-\Sigma_{\mu=1}^7I_\nu\mathscr{X}^ih_{t^\mu}I_\mu\mathscr{X}^i h\\&-\Sigma_{\mu=1}^7h_{t^\mu}I_\nu\mathscr{X}^iI_\mu\mathscr{X}^i h+\frac{1}{2}I_\nu\mathscr{X}^i f\mathscr{X}^i h+\frac{1}{2}fI_\nu\mathscr{X}^i\mathscr{X}^i h\bigg].
\end{aligned}\end{equation}
First, we should be careful here since the octonionic multiplication is not associative, we have
\begin{equation}
    \begin{aligned}
&h^{-1}\Sigma_{\nu=1}^7h_{t^\nu}\Sigma_{i=0}^7\bigg[2I_\nu\mathscr{X}^i h I_\nu\mathscr{X}^i h_{t^\nu}-\Sigma_{\mu=1}^7I_\nu\mathscr{X}^i h_{t^\mu}I_\mu\mathscr{X}^i h\bigg]\\=&h^{-1}\Sigma_{\nu=1}^7h_{t^{\nu}}\Sigma_{i=0}^7\bigg[I_\nu\mathscr{X}^i h I_\nu\mathscr{X}^i h_{t^\nu}-\Sigma_{\mu\neq \nu,\mu=1}^7I_\nu\mathscr{X}^i h_{t^\mu}I_\mu\mathscr{X}^i h \bigg]\\=&h^{-1}\Sigma_{\nu=1}^7h_{t^\nu}\Sigma_{i=0}^7 \Sigma_{\mu=1}^7I_{\mu}\mathscr{X}^i h_{t^\mu}I_\nu\mathscr{X}^i h\\=&\frac{1}{2}\Sigma_{\nu=1}^7h_{t^\nu}\Sigma_{i=0}^7\Sigma_{\mu=1}^7 B_\mu(\mathscr{X}^i)I_\nu\mathscr{X}^i h+\frac{7}{2}\Sigma_{\nu=1}^7h^{-2}h_{t^\nu}^2|\nabla_G h|^2.
 \end{aligned}
\end{equation}
The second equality holds by Lemma \ref{lemma2.2}, for $\nu\neq \mu$,
\begin{align*}
    I_\nu(-I_\nu I_\mu\mathscr{X}^i)=I_\mu \mathscr{X}^i, I_\mu(-I_\nu I_\mu\mathscr{X}^i)=-I_\nu\mathscr{X}^i,
\end{align*}
then we replace $\mathscr{X}^i$ by $-I_\nu I_\mu \mathscr{X}^i$ to get
\begin{align*}
\Sigma_{i=0}^7I_\nu\mathscr{X}^i h_{t^\mu}I_\mu\mathscr{X}^i h=\Sigma_{i=0}^7 I_\nu(-I_\nu I_\mu\mathscr{X}^i) h_{t^\mu}I_\mu(-I_\nu I_\mu\mathscr{X}^i) h=-\Sigma_{i=0}^7I_\mu \mathscr{X}^i h_{t^\mu}I_\nu \mathscr{X}^i h.
\end{align*}
The last equality is given by
\begin{align*}
    &\Sigma_{i=0}^7I_\mu\mathscr{X}^i h_{t^\mu}I_\nu\mathscr{X}^ih\\=&\Sigma_{i=0}^7\bigg[\frac{1}{2}hB_\mu(\mathscr{X}^i)I_\nu\mathscr{X}^i h+\frac{1}{2}\Sigma_{\tau=1}^7h^{-1}h_{t^\tau}I_\tau \mathscr{X}^i hI_\nu\mathscr{X}^i h-\frac{1}{4}h^{-1}f\mathscr{X}^i h I_\nu\mathscr{X}^i h\bigg]\\=&\Sigma_{i=0}^7\bigg[\frac{1}{2}hB_\mu(\mathscr{X}^i)I_\nu\mathscr{X}^i h+\frac{1}{2}h^{-1}h_{t^\nu}I_\nu \mathscr{X}^i hI_\nu\mathscr{X}^i h\bigg],\end{align*}
since from Lemma \ref{lemma2.3}, $\Sigma_{i=0}^7I_\nu \mathscr{X}^i h I_\mu \mathscr{X}^i h=0$, for $\nu\neq \mu$ and $\Sigma_{i=0}^7I_\nu\mathscr{X}^i h\mathscr{X}^i h=0.$


Next, we have
\begin{align*}
    \frac{1}{2}\Sigma_{i=0}^7fI_\nu\mathscr{X}^i\mathscr{X}^i h=8fh_{t^\nu},
\end{align*}
and
\begin{align*}
    2h\Delta_G h_{t^\nu}=48 h f_{t^\nu}=12h\Sigma_{i=0}^7B_\nu(\mathscr{X}^i)I_\nu\mathscr{X}^i h,\end{align*}
    since
\begin{align}\label{ft}
    f_{t^\nu}=-\frac{1}{4}h^{-2}h_{t^\nu}|\nabla_G h|^2+\frac{1}{2}h^{-1}\Sigma_{i=0}^7  I_\nu\mathscr{X}^i h I_\nu\mathscr{X}^i h_{t^\nu}=\frac{1}{4}\Sigma_{i=0}^7B_\nu(\mathscr{X}^i)I_\nu\mathscr{X}^i h.
    \end{align}
We also have
\begin{align*}
    -h^{-1}\Sigma_{i=0}^7\Sigma_{\nu,\mu=1}^7h_{t^\nu}h_{t^\mu}I_\nu \mathscr{X}^i I_\mu\mathscr{X}^i h=-h^{-1}\Sigma_{\nu=1}^7h^2_{t^\nu}\Delta_G h=h^{-1}\Sigma_{\nu=1}^7\bigg[-24fh_{t^\nu}^2+64h_{t^\nu}^2\bigg],
    \end{align*}
this is because for $\nu\neq \mu$, by Lemma \ref{lemma2.3}, $\Sigma_{i=0}^7(I_\nu \mathscr{X}^iI_\mu \mathscr{X}^i h+I_\mu\mathscr{X}^iI_\nu\mathscr{X}^ih)=0.$

    Using Lemma \ref{lemma3.2}, we have
\begin{align*}
   & \frac{1}{2}h^{-1}\Sigma_{\nu=1}^7h_{t^\nu}\Sigma_{i=0}^7I_\nu\mathscr{X}^i f\mathscr{X}^i h\\=&h^{-1}\Sigma_{\nu=1}^7h_{t^\nu}\bigg[\frac{1}{4}\Sigma_{i=0}^7hD(I_\nu\mathscr{X}^i)\mathscr{X}^i h +\frac{1}{2}h^{-1}h_{t^\nu}|\nabla_G h|^2\bigg]\\=&h^{-1}\Sigma_{\nu=1}^7h_{t^\nu}\bigg[-\frac{1}{4}\Sigma_{i=0}^7hD(\mathscr{X}^i)I_\nu\mathscr{X}^i h +\frac{1}{2}h^{-1}h_{t^\nu}|\nabla_G h|^2\bigg].\end{align*}
   Notice that
\begin{align*}
    64h^{-1}h^2_{t^\nu}+\frac{1}{2}h^{-2}h^2_{t^\nu}|\nabla_G h|^2=16h^{-1}h^2_{t^\nu}f-\frac{7}{2}h^{-2}h^2_{t^\nu}|\nabla_G h|^2,
    \end{align*}
we arrive at
\begin{align*}
   &h^{-1}\Sigma_{\nu=1}^7h_{t^\nu}\Sigma_{i=0}^7\bigg[2I_\nu\mathscr{X}^i h I_\nu\mathscr{X}^i h_{t^\nu}+2h I_\nu\mathscr{X}^i I_\nu\mathscr{X}^i h_{t^\nu}-\Sigma_{\mu=1}^7I_\nu\mathscr{X}^i h_{t^\mu}I_\mu\mathscr{X}^i_\alpha h\\&-\Sigma_{\mu=1}^7h_{t^\mu}I_\nu\mathscr{X}^i I_\mu\mathscr{X}^i h+\frac{1}{2}I_\nu\mathscr{X}^i f\mathscr{X}^i h+\frac{1}{2}fI_\nu\mathscr{X}^i\mathscr{X}^i h\bigg]\\=&\Sigma_{\nu=1}^7h_{t^\nu}\Sigma_{i=0}^7\bigg[-\frac{1}{4}D(\mathscr{X}^i )I_\nu\mathscr{X}^i h+12B_\nu(\mathscr{X}^i)I_\nu\mathscr{X}^ih+\frac{1}{2}\Sigma_{\mu=1}^7B_\mu(\mathscr{X}^i)I_\nu\mathscr{X}^i h\bigg].
    \end{align*}
Last, we have
\begin{align*}
    &\Sigma_{\nu=1}^7\Sigma_{i=0}^7hI_\nu\mathscr{X}^i h_{t^\nu}B_\nu(\mathscr{X}^i)\\=&\Sigma_{\nu=1}^7\Sigma_{i=0}^7hB_\nu(\mathscr{X}^i)\bigg[\frac{1}{2}hB_\nu(\mathscr{X}^i)+\frac{1}{2}\Sigma_{\mu=1}^7h^{-1}h_{t^\mu}(I_\mu\mathscr{X}^i)h-\frac{1}{4}h^{-1}f\mathscr{X}^i h\bigg]\\=&\frac{1}{2}\Sigma_{\nu=1}^7\Sigma_{i=0}^7h^2B_\nu(\mathscr{X}^i)^2+\Sigma_{i=0}^7\bigg[\frac{1}{2}\Sigma_{\mu=1}^7h_{t^\mu}(I_\mu\mathscr{X}^i)h B(\mathscr{X}^i)-\frac{1}{4}fB(\mathscr{X}^i )\mathscr{X}^i h\bigg].
\end{align*}
Combining these, we complete the proof.
\end{proof}

\begin{theorem}[Jerison-Lee type identity]We have

 

  \begin{equation}\label{jl}
     \begin{aligned}
        &\frac{1}{3}O_1+\frac{4}{3}O_2+4O_3\\=& \frac{1}{3}fh\mathcal{D}^2+\frac{1}{6}\Sigma_{i=0}^7h^2 D(\mathscr{X}^i)^2+
        \Sigma_{i=0}^7\Sigma_{\nu=1}^7h^2\bigg[\frac{2}{3}F_\nu(\mathscr{X}^i)B_\nu (\mathscr{X}^i)+2 B_\nu(\mathscr{X}^i)^2\bigg]\\\geq & \frac{1}{1000}\bigg[fh\mathcal{D}^2+\Sigma_{i=0}^7\Sigma_{\nu=1}^7h^2B_\nu (\mathscr{X}^i)^2\bigg]
     \end{aligned}
 \end{equation}\end{theorem}

\begin{proof}
   

From the definition (\ref{3.8}), we have
\begin{align}\label{*}
    \Sigma_{i,j,k=0}^7\mathbb{D}(\mathscr{X}^i,\mathscr{X}^j,\mathscr{X}^k)^2=16h^{-2}\mathcal{D}^2|\nabla h|^2+16\Sigma_{i=0}^7[D(\mathscr{X}^i)^2-\Sigma_{\nu=1}^7F_\nu(\mathscr{X}^i)^2]
\end{align}

This is because from Lemma \ref{lemma2.3}, we have
\begin{align*}
    \Sigma_{i,j,k=0}^7 \mathcal{D}(\mathscr{X}^i,\mathscr{X}^k)\mathscr{X}^j h\mathcal{D}(I_\nu\mathscr{X}^i,\mathscr{X}^k)I_\nu\mathscr{X}^j h=0\end{align*}
and 

\begin{align*}
    \Sigma_{i,j,k=0}^7 \mathcal{D}(I_\mu\mathscr{X}^i,\mathscr{X}^k)I_\mu\mathscr{X}^j h\mathcal{D}(I_\nu\mathscr{X}^i,\mathscr{X}^k)I_\nu\mathscr{X}^j h=0.\end{align*}

By the definition of $D(\mathscr{X}^i)$ and $F_\nu(\mathscr{X}^i)$, we have
\begin{align*}
    h^{-2}\Sigma_{i,j,k=0}^7\mathcal{D}(\mathscr{X}^i,\mathscr{X}^k)\mathscr{X}^j h \mathcal{D}(\mathscr{X}^j,\mathscr{X}^k)\mathscr{X}^i h=\Sigma_{i=0}^7D(\mathscr{X}^i)^2,\end{align*}
\begin{align*}
    &h^{-2}\Sigma_{i,j,k=0}^7\mathcal{D}(\mathscr{X}^i,\mathscr{X}^k)\mathscr{X}^j h \mathcal{D}(I_\nu\mathscr{X}^j,\mathscr{X}^k)I_\nu\mathscr{X}^i h\\=&-h^{-2}\Sigma_{i,j,k=0}^7\mathcal{D}(I_\nu\mathscr{X}^i,I_\nu\mathscr{X}^k)\mathscr{X}^j h \mathcal{D}(I_\nu\mathscr{X}^j,I_\nu\mathscr{X}^k)\mathscr{X}^i h\\=&-\Sigma_{i=0}^7F_\nu(\mathscr{X}^i)^2,\end{align*}
and
\begin{align*}
    &h^{-2}\Sigma_{i,j,k=0}^7\mathcal{D}(I_\nu \mathscr{X}^i,\mathscr{X}^k)I_\nu\mathscr{X}^j h \mathcal{D}(I_\nu \mathscr{X}^j,\mathscr{X}^k)I_\nu\mathscr{X}^i h\\=& h^{-2}\Sigma_{i,j,k=0}^7\mathcal{D}(\mathscr{X}^i,\mathscr{X}^k)\mathscr{X}^j h \mathcal{D}(\mathscr{X}^j,\mathscr{X}^k)\mathscr{X}^i h\\=&\Sigma_{i=0}^7D(\mathscr{X}^i)^2.\end{align*}
We should be careful to deal with the rest terms because of the non-associativity.
\begin{align*}
    &h^{-2}\Sigma_{\nu\neq \mu,\nu=1}^7\Sigma_{i,j,k=0}^7\mathcal{D}(I_\mu \mathscr{X}^i,\mathscr{X}^k)I_\mu \mathscr{X}^j h \mathcal{D}(I_\nu\mathscr{X}^j,\mathscr{X}^k)I_\nu\mathscr{X}^i h\\=&h^{-2}\Sigma_{\nu\neq \mu,\nu=1}^7\Sigma_{i,j,k=0}^7\mathcal{D}(I_\mu I_\nu \mathscr{X}^i,\mathscr{X}^k) \mathscr{X}^i h \mathcal{D}(I_\nu I_\mu \mathscr{X}^j,\mathscr{X}^k)\mathscr{X}^j h\\=  
    &-\Sigma_{\nu\neq \mu, \nu=1}^7\Sigma_{i=0}^7F_\nu(\mathscr{X}^i)^2,\end{align*}
since for $\nu=1,\mu=2$,
we have $I_1I_2\mathscr{X}^i=\delta I_3\mathscr{X}^i$ and $I_2I_1\mathscr{X}^i=-\delta I_3\mathscr{X}^i$ for $i=0,1,,...,7$, where $\delta=1$ or $-1$. Then
\begin{align*}
&h^{-2}\Sigma_{i,j,k=0}^7\mathcal{D}(I_2 I_1 \mathscr{X}^i,\mathscr{X}^k) \mathscr{X}^i h \mathcal{D}(I_1 I_2 \mathscr{X}^j,\mathscr{X}^k)\mathscr{X}^j h\\=& -\delta^2h^{-2}\Sigma_{i,j,k=0}^7\mathcal{D}(I_3\mathscr{X}^i,I_3\mathscr{X}^k)\mathscr{X}^i hD(I_3\mathscr{X}^j,I_3\mathscr{X}^k)\mathscr{X}^jh=-\Sigma_{i=0}^7F_3(\mathscr{X}^i)^2
\end{align*}






    To prove (\ref{jl}), we use (\ref{*}), that is
    \begin{align*}
        h^{-2}\mathcal{D}^2|\nabla h|^2+\Sigma_{i=0}^7[D(\mathscr{X}^i)^2-\Sigma_{\nu=1}^7F_\nu(\mathscr{X}^i)^2]\geq 0.    \end{align*}
        Then
\begin{align*}
        \frac{33}{100}fh\mathcal{D}^2+\frac{33}{400}\Sigma_{i=0}^7h^2D(\mathscr{X}^i)^2   \geq  \frac{33}{400}h^2\Sigma_{i=0}^7\Sigma_{\nu=1}^7F_\nu(\mathscr{X}^i)^2
    \end{align*}
 and 
 \begin{align*}
     \frac{33}{400}F_\nu(\mathscr{X}^i)^2+2B_\nu(\mathscr{X}^i)^2+\frac{2}{3}F_\nu(\mathscr{X}^i)B_\nu (\mathscr{X}^i)\geq \frac{1}{1000}B_\nu (\mathscr{X}^i)^2
 \end{align*}
 give (\ref{jl}).
\end{proof}



















\section{Proof of the main result }\label{s4}


\begin{thebibliography}{99}
\bibitem[CLMR2023]{CLMR2023} G. Catino, Y. Li, D. D. Monticell, A. Roncoroni. A Liouville Theorem in the Heisenberg group. To appear in J. Eur. Math. Soc. (JEMS)
\bibitem[CLZ]{CLZ} M. Christ; H. Liu; A. Zhang. Sharp Hardy-Littlewood-Sobolev inequalities on quaternionic Heisenberg groups. Nonlinear Anal. 130 (2016), 361–395.
\bibitem[CLZ2]{CLZ2} M. Christ; H. Liu; A. Zhang. Sharp Hardy-Littlewood-Sobolev inequalities on the octonionic Heisenberg group. Calc. Var. Partial Differential Equations 55 (2016), no. 1, Art. 11, 18 pp.
\bibitem[CMRW2024]{CMRW2024} G. Catino, D. D. Monticelli, A. Roncoroni, X. Wang. Liouville theorems on pseudohermitian manifolds with nonnegative Tanaka-Webster-Ricci curvature. arXiv:2412.08500
\bibitem[FL2012]{FL2012} R. L. Frank; E. H Lieb. Sharp constants in several inequalities on the Heisenberg group. Ann. of Math. (2) 176 (2012), no. 1, 349–381.
\bibitem[FS1974]{FS1974} G. B. Folland, E. M. Stein. Estimates for the $\bar{\partial}_b$ complex and analysis on the Heisenberg group. Comm. Pure Appl. Math. 27 (1974), 429-522.
\bibitem[FV2023]{FV2023} J. Flynn and J. Vétois, Liouville-type results for the CR Yamabe equation in the Heisenberg group. To appear in Ann. Sc. Norm. Super. Pisa Cl. Sci. (5)
\bibitem[G2001a]{G2001a} N. Gamara. The CR Yamabe conjecture$-$the case $n=1$. J. Eur. Math. Soc. (JEMS) 3 (2001), no. 2, 105–137.
\bibitem[G2001b]{G2001b} N. Gamara, R. Yacoub. CR Yamabe conjecture$-$the conformally flat case. Pacific J. Math. 201 (2001), no. 1, 121–175.
\bibitem[GV2001]{GV2001} N. Garofalo, D. Vassilev. Symmetry properties of positive entire solutions of Yamabe-type equations on groups of Heisenberg type. Duke Math. J. 106 (2001), no. 3, 411–448.
\bibitem[HW2022]{HW2022} F. Hang; X. Wang. A simpler proof of Frank and Lieb's sharp inequality on the Heisenberg group. arXiv:2211.10301
\bibitem[IMV2010]{IMV2010} S. Ivanov, I. Minchev, D. Vassilev. Extremals for the Sobolev inequality on the seven-dimensional quaternionic Heisenberg group and the quaternionic contact Yamabe problem. J. Eur. Math. Soc. (JEMS) 12 (2010), no. 4, 1041–1067.
\bibitem[IMV2014]{IMV2014} S. Ivanov, I. Minchev, D. Vassilev. Quaternionic contact Einstein structures and the quaternionic contact Yamabe problem. Mem. Amer. Math. Soc. 231 (2014), no. 1086, vi+82 pp.
\bibitem[IMV2023]{IMV2023} S. Ivanov, I. Minchev, D. Vassilev. Solution of the qc Yamabe equation on a 3-Sasakian manifold and the quaternionic Heisenberg group. Anal. PDE 16 (2023), no. 3, 839–860.
\bibitem[JL1987]{JL1987} D. Jerison, J. M. Lee. The Yamabe problem on CR manifolds. J. Differential Geom. 25 (1987), no. 2, 167–197.
\bibitem[JL1988]{JL1988} D. Jerison, J. M. Lee. Extremals for the Sobolev inequality on the Heisenberg group and the CR Yamabe problem. J. Amer. Math. Soc. 1 (1988), no. 1, 1–13.
\bibitem[JL1989]{JL1989} D. Jerison, J. M. Lee. Intrinsic CR normal coordinates and the CR Yamabe problem. J. Differential Geom. 29 (1989), no. 2, 303–343.
\bibitem[LZ]{LZ} D. Lin, Y. Zhou. Liouville type theorem for a semi-linear sub-elliptic equation on the quaternionic Heisenberg group. Preprint.
\bibitem[MO2023]{MO2023} X.-N. Ma, Q. Ou. A Liouville theorem for a class semilinear elliptic equations on the Heisenberg group. Adv. Math. 413 (2023), Paper No. 108851, 20 pp.
\bibitem[MOW2023]{MOW2023} X.-N. Ma, Q. Ou, T. Wu. Jerison-Lee identity and Semi-linear subelliptic equation on CR manifold. arXiv:2311.16428
\bibitem[W2015]{W2015} X. Wang. On a remarkable formula of Jerison and Lee in CR geometry. Math. Res. Lett. 22 (2015), no. 1, 279–299.
\bibitem[Y2024]{Y2024} Q. Yang, The optimal constant in the $L^2$ Folland-Stein inequality on the H-type group. J. Funct. Anal. 286 (2024), no. 2, Paper No. 110209, 22 pp.
\end{thebibliography}
\end{document}
